\documentclass[11pt,a4paper]{article}
\usepackage[T1]{fontenc}
\usepackage[utf8]{inputenc}
\usepackage{lmodern,microtype}
\usepackage[margin=26mm]{geometry}
\usepackage{amsmath,amssymb,amsthm,mathtools}
\usepackage{booktabs,tabularx,array}
\usepackage{xcolor}
\usepackage[colorlinks=true,linkcolor=blue!40!black,citecolor=blue!40!black,urlcolor=blue!40!black,breaklinks=true]{hyperref}
\usepackage{fancyhdr}
\pagestyle{fancy}\fancyhf{}
\fancyhead[L]{\small Integer-tail dynamics and slow modulations}
\fancyhead[R]{\small Research note / 28 October 2026}
\fancyfoot[C]{\thepage}
\setlength{\headheight}{14pt}
\setlength{\emergencystretch}{2em}
\newtheorem{theorem}{Theorem}[section]
\newtheorem{lemma}[theorem]{Lemma}
\newtheorem{proposition}[theorem]{Proposition}
\newtheorem{corollary}[theorem]{Corollary}
\theoremstyle{remark}\newtheorem{remark}[theorem]{Remark}
\newcommand{\Z}{\mathbb Z}\newcommand{\Q}{\mathbb Q}\newcommand{\R}{\mathbb R}\newcommand{\C}{\mathbb C}
\DeclareMathOperator{\cent}{cent}\DeclareMathOperator{\Log}{Log}
\title{Integer-Tail Dynamics and a P\'olya--Carlson\newline Exclusion for Slow Modulations}
\author{Research note prepared for Keio\\\small AI-assisted derivation; independent review pending}
\date{28 October 2026}
\hypersetup{pdftitle={Integer-Tail Dynamics and a Polya-Carlson Exclusion for Slow Modulations},pdfauthor={Research note prepared for Keio}}
\begin{document}
\maketitle
\begin{abstract}
We combine normalized integer reciprocal tails with analytic continuation of
coefficient-interpolated generating functions. The resulting criterion excludes
rational reciprocal sums for a class not treated by the preceding rational-profile
note: for every $\alpha>1$, $c\ne0$, $\kappa>0$, and $\eta>0$, positive integers with
\[
 b_n/\alpha^{2^n}=1+c/(\log(n+2))^\kappa+O(e^{-\eta n})
\]
have irrational reciprocal sum. Fractional-power profiles and nonzero periodic
logarithmic profiles are also covered. The key analytic inputs are the classical
Le Roy--Lindel\"of continuation theorem and the P\'olya--Carlson dichotomy.
We separately derive the exact centered-residue dynamics behind Erd\H{o}s
Problems 243 and 263, an elementary half-integer obstruction at an integer base,
and a finite-transient obstruction to purely local absorption arguments.
These results do not settle the unrestricted $o(1)$ perturbation in Problem 263.
Historical originality, external peer review, and proof-assistant verification
are not claimed.
\end{abstract}

\section{Target, results, and scope}
The first question in Erd\H{o}s Problem 263 asks whether
\begin{equation}\label{eq:263}
 b_n\in\Z_{>0},\qquad b_n/2^{2^n}\longrightarrow1
 \quad\Longrightarrow\quad \sum_n\frac1{b_n}\notin\Q.
\end{equation}
The current community database and the consulted formal-conjecture declaration
record this question as open \cite{database,formal263}. Its separate question about necessary growth for Type~2
irrationality sequences is not addressed here. Finitely many terms may always
be changed or deleted: their reciprocal sum is rational.

\begin{theorem}[Slow analytic modulations]\label{thm:main}
Let $\alpha>1$, $\nu\ge2$, $\eta>0$, and $c\in\R\setminus\{0\}$.
Let $(b_n)$ be positive integers, and suppose, for all sufficiently large $n$,
that either
\begin{align}
 \frac{b_n}{\alpha^{2^n}}&=1+\frac{c}{(\log(n+\nu))^\kappa}
                  +O(e^{-\eta n}),\qquad \kappa>0,\label{eq:logprofile}\\
 \text{or}\qquad
 \frac{b_n}{\alpha^{2^n}}&=1+\frac{c}{(n+\nu)^\theta}
                  +O(e^{-\eta n}),\qquad 0<\theta<1.\label{eq:powerprofile}
\end{align}
Then $\sum_n1/b_n$ is irrational. The errors need not be monotone, continuous,
or of one sign.
\end{theorem}

\begin{corollary}[Rounded examples]\label{cor:rounded}
For sufficiently large $N$, both series
\[
 \sum_{n=N}^{\infty}\frac1{\lfloor 2^{2^n}(1+1/\log(n+2))\rfloor},
 \qquad
 \sum_{n=N}^{\infty}\frac1{\lfloor 2^{2^n}(1+1/\sqrt{n+2})\rfloor}
\]
are irrational. The base $2$ may be replaced by any real $\alpha>1$.
Floors may be replaced by ceilings, nearest-integer rounding, or a bounded
integer change, provided the denominators remain positive.
\end{corollary}
\begin{proof}
The relative rounding error is $O(\alpha^{-2^n})$, hence
$O(e^{-\eta n})$ for every fixed $\eta>0$.
\end{proof}

\begin{theorem}[Periodic logarithmic modulations]\label{thm:periodic}
Let $t_n$ be a real periodic sequence, not identically zero, and let
$\alpha>1$, $\nu\ge2$, $\kappa>0$, $\eta>0$. If positive integers satisfy
\[
 \frac{b_n}{\alpha^{2^n}}
   =1+\frac{t_n}{(\log(n+\nu))^\kappa}+O(e^{-\eta n}),
\]
then $\sum_n1/b_n$ is irrational.
\end{theorem}
For example, $t_n=1+2(-1)^n$ permits a leading correction that changes sign.
No assertion is made for an arbitrary bounded, nonperiodic $t_n$.

\paragraph{Comparison with the preceding note.}
The previous rational-modulation manuscript allowed certain errors of
polynomial order, such as $O(n^{-5})$ around $1+1/n$ \cite{previous}.
The present note permits slower, nonrational leading profiles, but its
positive-growth arguments require an exponentially small relative remainder.
It is a complementary result, not a uniform improvement of the earlier error
bound. In particular, $1+1/\log n+o(1/\log n)$ is \emph{not} covered as a class.

\section{Exact integer tails and the modular reformulation}
The normalized tails provide the integer quantities used both in the new
analytic argument and in the known pseudo-greedy reformulation of
Problem~243 \cite{koizumi}.

\subsection{The integer certificate}
Fix $N$ and suppose $S_N=\sum_{n=N}^\infty1/b_n=p/q\in\Q$, $q>0$.
Write
\begin{equation}\label{eq:tails}
 T_n=\sum_{j=n}^\infty\frac1{b_j},\quad
 D_n=q\prod_{j=N}^{n-1}b_j,\quad C_n=D_nT_n.
\end{equation}
Then
\[
 C_n=p\prod_{j=N}^{n-1}b_j-
 q\sum_{j=N}^{n-1}\frac{\prod_{k=N}^{n-1}b_k}{b_j}
 \in\Z_{>0}.
\]
The fractions $C_n/D_n$ are deliberately \emph{not reduced}.
Their exact evolution is
\begin{equation}\label{eq:exactupdate}
 D_{n+1}=b_nD_n,\qquad C_{n+1}=b_nC_n-D_n.
\end{equation}
Conversely, integrality at $n=N$ with some integer $q>0$ already implies
rationality of $T_N$. Thus no irrationality criterion is hidden in the
normalization itself.

\subsection{Centered residues and pseudo-greedy dynamics}
Suppose $b_{n+1}/b_n^2\to1$. Quadratic growth gives
\[
 T_n=b_n^{-1}+b_{n+1}^{-1}+O(b_{n+1}^{-2}),\qquad
 T_n^{-1}=b_n-\frac{b_n^2}{b_{n+1}}+O(b_n^{-1}).
\]
Consequently
\begin{equation}\label{eq:gap}
 \varepsilon_n=T_n^{-1}+1-b_n\longrightarrow0.
\end{equation}
For sufficiently large $n$, $b_n$ is therefore the nearest integer to
$T_n^{-1}+1$. This also follows from Koizumi's stronger quantitative theorem
\cite[Theorem~1]{koizumi}.

For $C>0$ an integer, let $\cent_C(D)$ denote the unique integer congruent to
$D$ modulo $C$ in $[-C/2,C/2)$. Once $|\varepsilon_n|<1/2$, put
\[
 e_n=\cent_{C_n}(D_n)=C_n\varepsilon_n.
\]
Then the entire trajectory is deterministic:
\begin{equation}\label{eq:modular}
 \boxed{\quad
 b_n=\frac{D_n-e_n}{C_n}+1,\qquad
 C_{n+1}=C_n-e_n,\qquad D_{n+1}=b_nD_n.
 \quad}
\end{equation}
One further algebraic elimination yields
\begin{align}
 e_{n+1}&=\cent_{C_n-e_n}(e_nb_n^2),\label{eq:residue2}\\
 b_{n+1}&=b_n(b_n-1)+1+
       \frac{e_nb_n^2-e_{n+1}}{C_n-e_n}.\label{eq:secondorder}
\end{align}
For instance, modulo $C_{n+1}$ we have $C_n\equiv e_n$ and
$D_n=C_n(b_n-1)+e_n\equiv e_nb_n$; multiplying by $b_n$ proves
\eqref{eq:residue2}. The quotient in \eqref{eq:secondorder} is an integer.

\begin{proposition}[Absorption]\label{prop:absorb}
If $e_n=0$ at one step, then $e_j=0$ for every $j\ge n$, and
\[
 b_{j+1}=b_j^2-b_j+1\qquad(j\ge n).
\]
\end{proposition}
\begin{proof}
This follows immediately from \eqref{eq:residue2} and
\eqref{eq:secondorder}.
\end{proof}

For a positive rational initial tail the nearest-integer rule defines a
positive remainder at each step and these remainders decrease to zero.
Indeed, $b_n>T_n^{-1}$; if the remainders had a positive limit, the denominators
would be bounded, contradicting the repeated subtraction of their reciprocals.
Thus the rule really gives a reciprocal expansion of the initial rational.
Moreover, \eqref{eq:modular} gives
\[
 \frac{b_{n+1}}{b_n^2}
 =\frac{1-b_n^{-1}+\varepsilon_n/b_n}{1-\varepsilon_n}
       +\frac{1-\varepsilon_{n+1}}{b_n^2}.
\]
Hence $\varepsilon_n\to0$ implies $b_{n+1}/b_n^2\to1$.
This proves the following exact reformulation, also given by Koizumi:
\[
 \begin{gathered}
 \text{Problem 243 is true}\
 \Longleftrightarrow\
 \text{every rational pseudo-greedy trajectory satisfying }e_n/C_n\to0\\
           \text{eventually has }e_n=0.
 \end{gathered}
\]
It is not a proof of that statement. The current community database records
Problem~243 as open \cite{database}.

\subsection{The growth-base constraint cannot be discarded}
If $b_{n+1}/b_n^2\to1$ and $b_n\to\infty$, let
$\delta_n=\log b_{n+1}-2\log b_n=o(1)$. Then
\[
 \lambda=\lim_{n\to\infty}2^{-n}\log b_n
\]
exists, and
\begin{equation}\label{eq:basetail}
 \log b_n-2^n\lambda
   =-\sum_{j=n}^{\infty}2^{n-j-1}\delta_j=o(1).
\end{equation}
In particular $b_n\sim\alpha^{2^n}$ with $\alpha=e^\lambda>1$.
Positivity of $\lambda$ follows, for large $N$, from
$b_{n+1}\ge b_n^2/2$ and $b_N>2$.
A counterexample to \eqref{eq:263} must have the special value
$\lambda=\log2$. A nonabsorbing rational trajectory at some other growth base
would not automatically refute Problem~263.

\section{An elementary half-integer obstruction}
The implication from a positive solution of Problem~243 to the first question
of Problem~263 can be seen without invoking any transcendence theorem for
Sylvester growth constants.

\begin{lemma}[Half-integer asymptotics]\label{lem:half}
Suppose a positive integer sequence tends to infinity and eventually satisfies
$b_{n+1}=b_n^2-b_n+1$. There is $\alpha>1$ such that
\[
 b_n=\alpha^{2^n}+\frac12+O(\alpha^{-2^n}).
\]
Consequently it cannot be asymptotic to $a^{2^n}$ for an integer $a\ge2$.
\end{lemma}
\begin{proof}
Set $y_n=b_n-1/2$. Then exactly
\[
 y_{n+1}=y_n^2+\frac14,
 \qquad
 \log y_{n+1}-2\log y_n
       =\log\left(1+\frac1{4y_n^2}\right).
\]
With $\log\alpha=\lim2^{-n}\log y_n$, the same summation as in
\eqref{eq:basetail} gives
\[
 2^n\log\alpha-\log y_n
   =\sum_{j=n}^{\infty}2^{n-j-1}
         \log\left(1+\frac1{4y_j^2}\right)=O(y_n^{-2}).
\]
Here $y_{n+1}\ge y_n^2$ bounds the tail. Exponentiating yields
$\alpha^{2^n}=y_n+O(y_n^{-1})$, proving the asserted expansion.
If also $b_n\sim a^{2^n}$, the growth base must be $\alpha=a$.
The integer $b_n-a^{2^n}$ would then tend to $1/2$, an impossibility.
\end{proof}

\begin{corollary}
A positive solution of Problem~243 implies a positive solution of the first
question in Problem~263. This is an implication, not an equivalence.
\end{corollary}

The same reasoning supplies useful necessary conditions on a possible
counterexample to \eqref{eq:263}. In its eventual centered-residue trajectory,
$e_n\ne0$ for every $n$, since a zero residue would trigger the contradiction
in Lemma~\ref{lem:half}. Because $e_n/C_n\to0$ and $|e_n|\ge1$, necessarily
\begin{equation}\label{eq:necessaryC}
 C_n\longrightarrow\infty,\qquad \log C_n=o(n).
\end{equation}
The second assertion also follows from $C_{n+1}/C_n=1-\varepsilon_n\to1$.
Writing $u_n=b_n/2^{2^n}$, the product formula below shows in addition
\begin{equation}\label{eq:necessaryprod}
 \prod_{j=N}^{n-1}u_j\longrightarrow\infty.
\end{equation}
If the negative gap mass were finite,
$\sum_n\max(-\varepsilon_n,0)<\infty$, then
$C_{n+1}=C_n(1-\varepsilon_n)$ would give a uniform upper bound on $C_n$.
Thus this mass must diverge for a counterexample. These are restrictions,
not an exclusion of every possible counterexample.

\section{An analytic-continuation exclusion criterion}
The following criterion is the new mechanism used in this note. It separates
the arithmetic issue (integer tails) from the analytic one (continuation of
a comparison generating function).

\begin{lemma}[Tail product formula]\label{lem:product}
Let $A_n=\alpha^{2^n}$, $\alpha>1$, and $b_n=A_nu_n$ with $u_n\to1$.
Then
\[
 T_n=\frac1{b_n}\left(1+O(A_n^{-1})\right).
\]
Under the rationality assumption in \eqref{eq:tails},
\begin{equation}\label{eq:product}
 C_{N+m}=\frac q{A_N}
   \frac{\prod_{j=N}^{N+m-1}u_j}{u_{N+m}}
        \left(1+O(A_{N+m}^{-1})\right).
\end{equation}
\end{lemma}
\begin{proof}
Since $b_{n+1}\asymp b_n^2$, the terms after $b_n^{-1}$ sum to
$O(b_n^{-2})$. Also
$\prod_{j=N}^{N+m-1}A_j=A_{N+m}/A_N$. Substitution in
\eqref{eq:tails} proves \eqref{eq:product}.
\end{proof}

\begin{theorem}[Continuation versus integer growth]\label{thm:criterion}
Let $R_n>0$ tend to $1$, and put
\[
 H_m=\frac{\prod_{j=N}^{N+m-1}R_j}{R_{N+m}}\qquad(m\ge0).
\]
Assume
\begin{equation}\label{eq:Hgrowth}
 \log H_m=o(m),\qquad \frac{\log H_m}{\log m}\longrightarrow+\infty,
\end{equation}
and that $F_H(z)=\sum_{m\ge0}H_mz^m$ admits analytic continuation through
at least one point of the unit circle. If positive integers satisfy
\[
 b_n/\alpha^{2^n}=R_n+O(e^{-\eta n}),\qquad \alpha>1,\quad\eta>0,
\]
then $\sum_n1/b_n$ is irrational.
\end{theorem}
\begin{proof}
Suppose otherwise and use the integer tails $C_n$. Write
$b_n=A_nR_nv_n$, so $v_n=1+O(e^{-\eta n})$. Their product converges to a
positive limit $V_\infty$, and
\[
 \frac{\prod_{j=N}^{N+m-1}v_j}{v_{N+m}}
        =V_\infty(1+O(e^{-\eta m})).
\]
By Lemma~\ref{lem:product}, for a constant $K>0$,
\begin{equation}\label{eq:CHe}
 C_{N+m}=KH_m+E_m,\qquad E_m=O(e^{-\eta m/2}).
\end{equation}
The exponential bound uses $H_m=e^{o(m)}$ and the doubly exponential
smallness of $A_{N+m}^{-1}$. Thus
\[
 F_C(z):=\sum_{m\ge0}C_{N+m}z^m=K F_H(z)+F_E(z),
\]
where $F_E$ is holomorphic in a disk of radius greater than $1$.
The integer-coefficient series $F_C$ has radius of convergence exactly $1$
and extends through the same boundary point as $F_H$.

The P\'olya--Carlson theorem says that an integer-coefficient series of
radius $1$ is rational or has the whole unit circle as a natural boundary;
see \cite[Theorem~1.1]{pc}. Hence $F_C$ is rational. Partial fractions of a
rational function with no poles inside the unit disk imply
$C_{N+m}=O(m^M)$ for some fixed $M$. This contradicts
\eqref{eq:Hgrowth} and \eqref{eq:CHe}.
\end{proof}

\begin{remark}[A complementary elementary case]\label{rem:decay}
With the same summable multiplicative error, if instead $H_m\to0$, then
$C_{N+m}\sim KH_m\to0$, directly contradicting positive integrality.
No continuation theorem is needed in that case.
\end{remark}

\section{Building the analytic continuation explicitly}
We now verify, rather than assume, the analytic hypothesis for the profiles in
Theorem~\ref{thm:main}.

\subsection{The classical continuation input}
We use the following form of the Le Roy--Lindel\"of theorem
\cite[Introduction]{leroy}: if $\Phi$ is holomorphic in a neighborhood of the
closed right half-plane and, for every $\epsilon>0$,
\[
 |\Phi(z)|\le C_\epsilon e^{\epsilon|z|}\qquad(\Re z\ge0),
\]
then $\sum_{m\ge0}\Phi(m)w^m$ continues off the nonnegative real axis.
Combined with its original unit disk, this supplies continuation to the slit
plane $\C\setminus[1,\infty)$ when the radius is $1$.
This is the zero-exponential-type case of the classical theorem.
It can also be obtained by the Lindel\"of contour representation: the factor
$\pi/\sin(\pi z)$ gives vertical exponential decay and the zero-type bound
allows the residue contour to be shifted to the right.

\subsection{An interpolation by convergent differences}
For either profile, take $N$ so large that all branches below are holomorphic
for $\Re z>N-1$, and define
\[
 f(z)=\Log R(z),
\]
where
\[
 R(z)=1+\frac{c}{(\Log(z+\nu))^\kappa}
 \quad\text{or}\quad
 R(z)=1+\frac{c}{(z+\nu)^\theta}.
\]
All logarithms and powers are the principal branches on that half-plane;
by increasing $N$ we ensure $|R(z)-1|<1/2$ there. Define
\begin{equation}\label{eq:interpolant}
 g(z)=\sum_{k=0}^{\infty}\bigl(f(N+k)-f(N+k+z)\bigr),\qquad
 \Phi(z)=\exp\bigl(g(z)-f(N+z)\bigr).
\end{equation}
The series is locally uniformly convergent for $\Re z>-1/2$, because the
following derivative bounds are summable along the positive integers:
\begin{equation}\label{eq:fbounds}
 \begin{array}{c|c|c}
 & |f(w)| & |f'(w)|\\ \hline
 \text{logarithmic} & O((\log|w+\nu|)^{-\kappa}) &
  O\bigl(|w+\nu|^{-1}(\log|w+\nu|)^{-\kappa-1}\bigr)\\[2pt]
 \text{fractional power} & O(|w+\nu|^{-\theta}) &
  O(|w+\nu|^{-1-\theta})
 \end{array}
\end{equation}
For an integer $m\ge0$, telescoping finite sums in \eqref{eq:interpolant}
and letting the upper limit tend to infinity gives
\[
 g(m)=\sum_{j=N}^{N+m-1}f(j),\qquad \Phi(m)=H_m.
\]
The omitted last $m$ terms tend to zero for each fixed $m$.

It remains to check the growth in the entire right half-plane, not just on
the positive axis. Write $r=|z|\ge2$, $\Re z\ge-1/4$, and split the sum
at $k=2r$. For $k\le2r$, use the two separate bounds on $f$ in
\eqref{eq:fbounds}; for $k>2r$, integrate $f'$ on the segment from
$N+k$ to $N+k+z$. Uniformly in this half-plane one obtains
\begin{align*}
 |g(z)|&\le C\left(1+
   \sum_{k\le2r}\frac1{(\log(N+k))^\kappa}
   +r\sum_{k>2r}\frac1{(N+k)(\log(N+k))^{\kappa+1}}\right)\\
 &=O\left(\frac r{(\log r)^\kappa}\right)
 \qquad\text{in the logarithmic case},\\
 |g(z)|&\le C\left(1+\sum_{k\le2r}(N+k)^{-\theta}
       +r\sum_{k>2r}(N+k)^{-1-\theta}\right)\\
 &=O(r^{1-\theta})\qquad\text{in the power case}.
\end{align*}
The fixed shifts in these bounds are harmless. The same estimates, with
an extra bounded term, bound $\log|\Phi(z)|$ from above. Both bounds are
$o(r)$. Thus $\Phi$ has zero exponential type on a neighborhood of the
closed right half-plane, as required. The generating function $F_H$
therefore has the desired continuation.

\subsection{Growth and proof of the main theorem}
For $c>0$, summing $f(n)$ gives
\begin{align}
 \log H_m&=(c+o(1))\frac m{(\log m)^\kappa}
                  &&\text{under \eqref{eq:logprofile}},\label{eq:loggrowth}\\
 \log H_m&=\left(\frac c{1-\theta}+o(1)\right)m^{1-\theta}
                  &&\text{under \eqref{eq:powerprofile}}.\label{eq:powergrowth}
\end{align}
For example, $\log(1+x)=x+O(x^2)$ and the quadratic correction is
$o(\sum x_n)$ for each of these slowly decaying positive sequences.
The elementary estimates
$\sum_{j\le m}(\log j)^{-\kappa}\sim m/(\log m)^\kappa$ and
$\sum_{j\le m}j^{-\theta}\sim m^{1-\theta}/(1-\theta)$ complete the calculation.
Both expressions satisfy \eqref{eq:Hgrowth}. Theorem~\ref{thm:criterion}
now proves irrationality. For $c<0$, the same asymptotics give
$\log H_m\to-\infty$, and Remark~\ref{rem:decay} applies. This proves
Theorem~\ref{thm:main}.

\section{Periodic profiles: sign changes are allowed}
Here we prove Theorem~\ref{thm:periodic}. Let $p$ be a period and
\[
 \mu=\frac1p\sum_{r=0}^{p-1}t_{N+r},\qquad
 v=\frac1p\sum_{r=0}^{p-1}t_{N+r}^2.
\]
If $\mu\ne0$, periodic summation and the expansion of $\log(1+x)$ give
\[
 \log H_m=(\mu+o(1))m/(\log m)^\kappa.
\]
Thus $\mu<0$ is handled by decay, while $\mu>0$ has the required
superpolynomial, subexponential growth.

For $\mu>0$ one must interpolate residue classes separately; replacing
$t_n$ by a trigonometric function of a complex argument would not provide
the required growth bound. Define, for $0\le r<p$,
\[
 f_r(z)=\Log\left(1+
 \frac{t_{N+r}}{(\Log(N+r+pz+\nu))^\kappa}\right),
 \qquad F(z)=\sum_{r=0}^{p-1}f_r(z).
\]
Increasing $N$ again makes all these functions holomorphic for
$\Re z>-1/2$. Set
\[
 G(z)=\sum_{k=0}^{\infty}(F(k)-F(k+z)),\qquad
 \Phi_r(z)=\exp\left(G(z)+\sum_{j=0}^{r-1}f_j(z)-f_r(z)\right).
\]
Then $\Phi_r(k)=H_{pk+r}$. The same prefix/tail estimates as above prove
zero exponential type for every $\Phi_r$. Hence each
$\sum_{k\ge0}H_{pk+r}w^k$ continues through every point of the unit circle
except possibly $1$. The decomposition
\[
 F_H(z)=\sum_{r=0}^{p-1}z^r\sum_{k\ge0}H_{pk+r}(z^p)^k
\]
gives continuation through every $z$ with $|z|=1$ and $z^p\ne1$.
Theorem~\ref{thm:criterion} applies.

Finally suppose $\mu=0$. Because $t_n$ has bounded partial sums,
$\sum t_n/(\log(n+\nu))^\kappa$ converges by summation by parts.
Since $t_n$ is not identically zero, $v>0$, and the quadratic term in
$\log(1+x)$ gives
\[
 \log H_m=-\left(\frac v2+o(1)\right)
                  \frac m{(\log m)^{2\kappa}}\longrightarrow-\infty.
\]
The cubic remainder is $o(m/(\log m)^{2\kappa})$.
Remark~\ref{rem:decay} completes the proof.

\section{Why finite local tests cannot close the absorption gap}
The following construction illustrates a genuine obstacle. Related
flat-residue constructions are presented in the independent working report
\cite{flat}; the proof here is explicit and does not depend on an unverified
claim from that report.

\begin{proposition}[Arbitrarily long small nonzero gaps]\label{prop:flat}
Given $L\ge1$ and $\epsilon>0$, a positive rational initial tail can be
chosen so that its centered-residue expansion has $L$ consecutive nonzero
gaps of magnitude less than $\epsilon$, followed by a zero gap and absorption.
\end{proposition}
\begin{proof}
Choose a prime $p>L+\max(3,\epsilon^{-1})$ and put
\[
 C=p-L,\qquad M=\prod_{j=0}^{L-1}(C+j).
\]
Since $p\nmid M$, choose $T>0$ with $MT\equiv1\pmod p$.
Increasing $T$ by a multiple of $p$, arrange $MT-1>10C$.
Start the unreduced trajectory at $C_0=C$, $D_0=MT-1$.

Let $M_j=\prod_{i=j}^{L-1}(C+i)$. We show inductively that for $j<L$,
\[
 C_j=C+j,\quad D_j\equiv-1\pmod{M_j},\quad p\mid D_j.
\]
These assertions hold at $j=0$. Since $C_j\ge3$, the centered residue is
$e_j=-1$. Therefore $C_{j+1}=C_j+1$. Also
$b_j=(D_j+1)/C_j+1\equiv1\pmod{M_{j+1}}$, whence
$D_{j+1}=D_jb_j\equiv-1\pmod{M_{j+1}}$. Divisibility by $p$ persists.
At step $L$, $C_L=p\mid D_L$, so $e_L=0$ and absorption follows.
The nonzero gaps are exactly $-1/(C+j)$, all of magnitude less than
$\epsilon$. The recurrence and the final telescoping Sylvester tail show
that the reciprocal sum is the initial rational $C_0/D_0$.
\end{proof}

This is not a counterexample to either problem: the gaps eventually vanish,
and the limiting growth base is not an integer by Lemma~\ref{lem:half}.
It does rule out a uniform finite-window argument asserting that a fixed
number of sufficiently small gaps forces an immediate zero gap. Treating
centered residues as independent random samples would likewise be a
heuristic, not a proof.

\section{The remaining full-problem obstruction}
For an unrestricted perturbation $b_n/2^{2^n}=1+o(1)$, the exact conditions
proved above imply that any rational reciprocal sum would produce a
rational initial state and a deterministic trajectory with
\begin{equation}\label{eq:openbridge}
 \boxed{\qquad
 e_n\ne0,\quad e_n/C_n\longrightarrow0,\quad
 \lim_{n\to\infty}2^{-n}\log b_n=\log2.
 \qquad}
\end{equation}
Conversely, a rational pseudo-greedy trajectory satisfying these conditions
would refute the first question in Problem~263, by \eqref{eq:basetail}.
Thus exclusion of these trajectories is an exact remaining task.

The P\'olya--Carlson theorem does not by itself exclude them. Their
integer-tail generating function has radius $1$, but may have a natural
boundary. No analytic continuation follows merely from
$C_{n+1}/C_n\to1$. Nor have we proved that the rational-generating-function
branch forces absorption in every case. In our explicit positive profiles,
continuation excludes the natural-boundary branch and superpolynomial
growth excludes the rational branch; both additional ingredients matter.

Similarly, the usual small-positive-integer argument at the unmodified tail
level stalls at this critical growth scale:
\[
 \left(\prod_{j=N}^{n-1}b_j\right)T_n
 \asymp \prod_{j=N}^{n-1}(b_j/2^{2^j}).
\]
The right side can grow rather than tend to zero. Fixed-order finite
differences alone also do not make a profile of size
$\exp(n/(\log n)^\kappa)$ tend to zero. The analytic continuation step is
precisely what replaces that earlier approach, not a proof that every
possible irregular profile is covered.

\section{Provenance and verification record}
Koizumi's published paper establishes the pseudo-greedy framework and proves
that all but countably many real bases are Type~2 irrationality bases
\cite{koizumi}. It does not identify base $2$ as a nonexceptional base.
Duverney's survey includes the classical integer-base criterion requiring
summable relative perturbations \cite[Theorem~2.3]{duverney}.
The leading positive perturbations $1/(\log n)^\kappa$ and $n^{-\theta}$,
$0<\theta<1$, are not summable. Consequently that particular criterion
cannot simply be substituted for the present proof. This comparison is
not an exhaustive priority search.

The complex-analytic ingredients are classical and explicitly cited.
The present note proves their hypotheses for the displayed profiles and
derives the resulting exclusion criterion. It makes no claim that this
application is historically first. No unproved result in the user's corpus
or in the preceding research note is used as a premise.

The accompanying \texttt{verify\_exact.py} uses only integer and rational
arithmetic. It checked 10,000 centered-residue states and their two-step
identities, the trajectory beginning at $11/29$, flat transients of lengths
$1$ through $10$, eight exact rounded square-root-profile denominators, and
1,000 instances of the half-shift identity. All checks passed in the saved
run. These finite checks do not verify the complex-analytic arguments or
prove an infinite irrationality theorem. No Lean or other proof-assistant
verification, external peer review, or unrestricted solution of Problem~263
is asserted.

\begin{thebibliography}{9}
\small
\bibitem{database}
T. Tao and contributors, \emph{Erd\H{o}s problem database},
entries 243 and 263, consulted 28 October 2026.
\url{https://github.com/teorth/erdosproblems}.

\bibitem{formal263}
The Formal Conjectures Authors, \emph{Erd\H{o}s Problem 263},
\path{FormalConjectures/ErdosProblems/263.lean}, consulted 28 October 2026.
\href{https://github.com/google-deepmind/formal-conjectures/blob/main/FormalConjectures/ErdosProblems/263.lean}{Formal Conjectures repository}.
The main question is tagged \texttt{research open}; its \texttt{sorry}
is a placeholder, not a proof.

\bibitem{koizumi}
J.~Koizumi, \emph{Irrationality of the reciprocal sum of doubly exponential
sequences}, Integers \textbf{26} (2026), A28.
\url{https://math.colgate.edu/~integers/aa28/aa28.pdf}.
Also arXiv:2504.05933.

\bibitem{pc}
J.~P. Bell, K. Gunn, K.~D. Nguyen, and J.~C. Saunders,
\emph{A general criterion for the P\'olya--Carlson dichotomy and application},
arXiv:2206.00862 (2022), Theorem~1.1.
\url{https://arxiv.org/html/2206.00862}.

\bibitem{leroy}
A.~Mkrtchyan, \emph{Multivariate analog of the Le Roy--Lindel\"of theorem
about power series analytic continuation}, arXiv:2403.17021 (2024),
Introduction, statement of the classical one-variable theorem.
\url{https://arxiv.org/html/2403.17021}.

\bibitem{duverney}
D.~Duverney, \emph{Irrationality of fast converging series of rational
numbers}, J. Math. Sci. Univ. Tokyo \textbf{8} (2001), 275--316.
\url{https://www.ms.u-tokyo.ac.jp/journal/pdf/jms080206.pdf}.

\bibitem{flat}
\emph{Erd\H{o}s Problem a Day: Problem 243}, independent working report,
dated 12 August 2026, consulted 28 October 2026.
\url{https://www.erdosproblemaday.com/report/243}.
Cited for prior flat-transient constructions, not as an independently
refereed resolution of either problem.

\bibitem{previous}
\emph{A Rational-Modulation Exclusion for Doubly Exponential Reciprocal
Sums}, research note prepared for Keio, 28 October 2026.
Earlier manuscript in the present research session; not a proof dependency.
\end{thebibliography}
\end{document}
